\bookchapter{Κεφάλαιο 55: Έργο 2 — Εφαρ$\mu$ογή αναζήτησης καιρού (δίκτυο)}

\chapterintro{Το βιβλίο εξόδων ήταν «εσωτερική εξάσκηση»· η εφαρ$\mu$ογή καιρού είναι «έξοδος στον κόσ$\mu$ο» — έχει να κάνει $\mu$ε το διαδίκτυο. Σε αυτό το έργο θα $\mu$άθεις: να κάνεις αίτηση για δωρεάν API, να αναλύεις JSON, να διαχειρίζεσαι τις τρεις καταστάσεις διεπαφής «φόρτωση/επιτυχία/αποτυχία», και πώς να κρύβεις καλά το κλειδί API. Όταν το τελειώσεις, θα έχεις την ικανότητα να φτιάχνεις οποιαδήποτε «συνδεδε$\mu$ένη εφαρ$\mu$ογή» — ειδήσεις, ισοτι$\mu$ίες, $\mu$ετάφραση: ίδια συνταγή.}

\booksection{Αίτηση για δωρεάν API καιρού}

Για προσωπική εξάσκηση προτείνονται το OpenWeatherMap ή το QWeather (δη$\mu$οφιλές στην Κίνα) — και τα δύο έχουν δωρεάν όριο, αρκεί να γραφτείς. Με παράδειγ$\mu$α το OpenWeatherMap:

\begin{enumerate}
  \item Δη$\mu$ιουργία λογαριασ$\mu$ού στη σελίδα API keys δη$\mu$ιουργείς ένα κλειδί (apikey)·
  \item Η δωρεάν έκδοση $\mu$πορεί να καλέσει το Current Weather (τρέχουσα πρόγνωση): api.openweathermap.org/data/2.5/weather? q=beijing\&appid=ΤΟ\_ΚΛΕΙΔΙ\_ΣΟΥ\&units=metric·
  \item Δοκί$\mu$ασέ το $\mu$ια φορά στον browser και δες πώς $\mu$οιάζει το JSON που επιστρέφεται — αυτό είναι το υποχρεωτικό «διάβασ$\mu$α» πριν συνδεθείς $\mu$ε οποιοδήποτε API.
\end{enumerate}

\textbf{Προσοχή στην παγίδα: το κλειδί δεν πρέπει ποτέ να γραφτεί σκληρά στον κώδικα}\par

Το API Key είναι σαν τον «κωδικό του πορτοφολιού σου» — $\mu$όλις $\mu$πει σε κώδικα και γίνει commit στο GitHub, $\mu$έσα σε λίγες ώρες θα το αρπάξουν crawlers για κατάχρηση. Η σωστή $\mu$έθοδος: βάλ’ το στο local.properties (αρχείο που το gitignore αγνοεί ήδη) και διάβασέ το στον κώδικα $\mu$έσω του buildConfig του Gradle. Ψάξε «best practices για την απόκρυψη του API Key σε Android» και ακολούθησέ το.

\booksection{Στρώ$\mu$α δικτύου: Retrofit + ανάλυση JSON}

Ορίσε πρώτα το «$\mu$οντέλο δεδο$\mu$ένων» σύ$\mu$φωνα $\mu$ε το JSON που επιστρέφει το API· όσα πεδία δεν ταιριάζουν, αγνόησέ τα:

\begin{lstlisting}[language=Kotlin]
// JSON: {"weather":[{"description":"Αίθριο"}],"main":{"temp":26.5,"humidity":40},"name":"Πεκίνο"}
data class WeatherResponse(
    val weather: List<Weather>,
    val main: Main,
    val name: String
)
data class Weather(val description: String)
data class Main(val temp: Double, val humidity: Int)

interface WeatherApi {
    @GET("data/2.5/weather")
    suspend fun getWeather(
        @Query("q") city: String,
        @Query("appid") key: String,
        @Query("units") metric: String = "metric",
        @Query("lang") lang: String = "zh_cn"
    ): WeatherResponse
}
\end{lstlisting}

Το Gson θα στρώσει αυτό$\mu$ατα το JSON σε αυτά τα data class — αρκεί τα ονό$\mu$ατα των πεδίων να ταιριάζουν και η $\mu$αγεία γίνεται $\mu$όνη της.

\booksection{Τρεις καταστάσεις διεπαφής: φόρτωση / επιτυχία / αποτυχία}

Η ψυχή $\mu$ιας συνδεδε$\mu$ένης εφαρ$\mu$ογής είναι η «διαχείριση καταστάσεων». Αργό δίκτυο, αποσύνδεση, λάθος όνο$\mu$α πόλης… για κάθε περίπτωση η διεπαφή πρέπει να δίνει απάντηση στον χρήστη. Εκφράσ’ το κο$\mu$ψά $\mu$ε ένα sealed class (σφραγισ$\mu$ένη κλάση):

\begin{lstlisting}[language=Kotlin]
sealed class UiState {
    object Loading : UiState()                    // περιστρέφεται
    data class Success(val data: WeatherResponse) : UiState()  // επιτυχία
    data class Error(val message: String) : UiState()          // αποτυχία
}

// Στο ViewModel:
fun search(city: String) {
    viewModelScope.launch {
        _state.value = UiState.Loading
        _state.value = try {
            UiState.Success(api.getWeather(city))
        } catch (e: Exception) {
            UiState.Error("Δεν βρήκα αυτή την πόλη, δοκί$\mu$ασε άλλο όνο$\mu$α;")
        }
    }
}

// Στην οθόνη:
when (val s = state) {
    is UiState.Loading -> CircularProgressIndicator()   // περιστροφή
    is UiState.Success -> WeatherContent(s.data)        // ε$\mu$φάνιση
    is UiState.Error -> Text(s.message)                 // ε$\mu$φάνιση σφάλ$\mu$ατος
}
\end{lstlisting}

Το τρι$\mu$ερές Loading/Success/Error είναι το «πανάκεια πρότυπο» των συνδεδε$\mu$ένων διεπαφών — $\mu$άθε το απ’ έξω και $\mu$ετά, σε όλες τις συνδεδε$\mu$ένες εφαρ$\mu$ογές, ακολούθησε το ίδιο σχή$\mu$α.

\booksection{Κώδικας διεπαφής στην πράξη: να ζωγραφίσεις τον καιρό}

Με τις καταστάσεις έτοι$\mu$ες, η διεπαφή $\mu$οιάζει έτσι (κρίσι$\mu$α αποσπάσ$\mu$ατα):

\begin{lstlisting}[language=Kotlin]
@Composable
fun WeatherContent(data: WeatherResponse) {
    Column(
        Modifier.fillMaxSize().padding(24.dp),
        horizontalAlignment = Alignment.CenterHorizontally
    ) {
        Text(data.name, fontSize = 28.sp, fontWeight = FontWeight.Bold)
        Text(
            "${data.main.temp}°C",
            fontSize = 72.sp,
            fontWeight = FontWeight.Light
        )
        Text(
            data.weather.firstOrNull()?.description ?: "Άγνωστο",
            fontSize = 20.sp,
            color = Color.Gray
        )
        Spacer(Modifier.height(16.dp))
        Row {
            Text(" Υγρασία ${data.main.humidity}%")
            Spacer(Modifier.width(24.dp))
            Text(" Αίσθηση ${data.main.temp}°C")
        }
    }
}
\end{lstlisting}

Μερικές λεπτο$\mu$έρειες αξίζουν προσοχή: το firstOrNull() παίρνει $\mu$ε ασφάλεια το πρώτο στοιχείο (η λίστα $\mu$πορεί να είναι κενή), και το ?: “Άγνωστο” δίνει κεί$\mu$ενο εφεδρείας — η αρχή null-safety ισχύει και στη διεπαφή. Η θερ$\mu$οκρασία $\mu$ε $\mu$εγάλο $\mu$έγεθος γρα$\mu$$\mu$ατοσειράς 72.sp για εντύπωση και το ελαφρύ βάρος (Light) κάνουν τους αριθ$\mu$ούς κο$\mu$ψούς — κλασική συνταγή εφαρ$\mu$ογών καιρού.

\booksection{Λεπτο$\mu$έρειες διεπαφής και αλληλεπίδρασης}

\booksection{Λεπτο$\mu$έρειες διεπαφής και αλληλεπίδρασης}

\begin{itemize}
  \item Πεδίο αναζήτησης: TextField στην κορυφή, γράφεις όνο$\mu$α πόλης και πατάς αναζήτηση ($\mu$ε ενέργεια «αναζήτηση» στο πληκτρολόγιο — Enter και ψάχνει)·
  \item Κάρτα καιρού: θερ$\mu$οκρασία σε $\mu$εγάλα γρά$\mu$$\mu$ατα (26.5°), περιγραφή καιρού (ηλιοφάνεια), υγρασία, όνο$\mu$α πόλης·
  \item Αγαπη$\mu$ένες πόλεις: αποθήκευσε τις πόλεις που κοιτάς συχνά στο DataStore, ώστε να αναζητούνται αυτό$\mu$ατα στο επό$\mu$ενο άνοιγ$\mu$α — και $\mu$ια επανάληψη του κεφαλαίου 6 κάτω από το χέρι·
  \item Τράβηγ$\mu$α για ανανέωση: στο Compose $\mu$ε το PullToRefreshBox (το δίνει το Material 3), και η αίσθηση γίνεται α$\mu$έσως «επαγγελ$\mu$ατική».
\end{itemize}

\begin{tcbinfo}{Συ$\mu$βουλή: οι εικόνες ζωντανεύουν τη διεπαφή}
Για κάθε είδος καιρού βάλε ένα emoji ή εικονίδιο (), $\mu$ε αντιστοίχιση when(description) — τρεις γρα$\mu$$\mu$ές κώδικα και η «θερ$\mu$οκρασία» της διεπαφής διπλασιάζεται. Μπορείς και να φορτώσεις βιβλιοθήκη εικονιδίων καιρού τρίτων, αλλά στο στάδιο της εξάσκησης τα emoji αρκούν.

\end{tcbinfo}

\begin{tcbpractice}{Πρακτική άσκηση: λίστα ελέγχου της εφαρ$\mu$ογής καιρού}
(1) Η αναζήτηση καιρού δουλεύει κανονικά· (2) $\mu$ε αποσύνδεση δικτύου, η εφαρ$\mu$ογή ε$\mu$φανίζει κο$\mu$ψό σφάλ$\mu$α χωρίς κατάρρευση (δοκί$\mu$ασέ το σε λειτουργία πτήσης!)· (3) $\mu$ε ανύπαρκτη πόλη, φιλικό $\mu$ήνυ$\mu$α· (4) $\mu$ε δύο αγαπη$\mu$ένες πόλεις, να $\mu$ένουν $\mu$ετά από επανεκκίνηση της εφαρ$\mu$ογής· (5) $\mu$ε περιστροφή οθόνης να $\mu$η χάνεται η κατάσταση. Αν περάσουν και οι πέντε, το έργο «αποφοιτά».

\end{tcbpractice}

\begin{tcbsummary}{Σύνοψη κεφαλαίου}
\begin{itemize}
  \item Πρώτα δοκι$\mu$ή του API στον browser και κατανόηση του JSON, $\mu$ετά κώδικας — το «προπαρασκευαστικό διάβασ$\mu$α» δεν χάνει ποτέ
  \item Το API Key κρύβεται στο local.properties, ποτέ $\mu$έσα στη βάση κώδικα
  \item Interface Retrofit + data class: $\mu$ε τα σωστά πεδία, το JSON αναλύεται αυτό$\mu$ατα
  \item Η τρι$\mu$ερής κατάσταση Loading/Success/Error είναι το πανάκεια των συνδεδε$\mu$ένων διεπαφών
  \item Αποσύνδεση, κενή είσοδος, ανύπαρκτη πόλη — η διαχείριση σφαλ$\mu$άτων είναι η βάση επιτυχίας κάθε συνδεδε$\mu$ένης εφαρ$\mu$ογής
\end{itemize}

\end{tcbsummary}
